\honest{Magical Categories}{Eliezer Yudkowsky}{2008}

I open with Bill Hibbard's proposal of 2001: build simple machines that ``learn to recognize happiness and unhappiness in human facial expressions, human voices and human body language,'' then hard-wire the result as the values of more intelligent machines. It was ``published in a peer-reviewed journal,'' and Hibbard later wrote a book about it, so it is not a strawman. \nb{Hibbard's own page calls the piece a column, in \textsc{Computer Graphics}; peer review could not be confirmed. The column also proposed predicting future happiness and measures such as health and wealth.}

When I asked whether such an AI would tile the galaxy with molecular smiley faces, Hibbard replied that a superintelligence's hard-wired recognition ``will certainly not be fooled'' by them. He also wrote of my ``preference for drama over reason,'' which leads me to remark that most would-be AI builders are unsuited to the task. That is not my topic today.

Once upon a time, the US Army trained a neural network on 50 photos of camouflaged tanks among trees and 50 of trees alone; held back from 200, the other 100 were all classified correctly. The Pentagon's own tests did no better than chance: the tank photos had been taken on cloudy days and the forest photos on sunny ones, and the network had learned the weather. I have never found an original source, and the story ``might or might not be fact.'' \nb{Gwern Branwen's investigation found it to be an urban legend, probably grown from a question Edward Fredkin asked in the 1960s.} The lesson: if training cases and real cases come from different processes, past success guarantees nothing. The problem is not unsolvable, but it is unpatchable; there are deep answers but no bandaids.

With 256 by 256 black-and-white images there are $2^{65536}$ images and $2^{2^{65536}}$ ways to divide them, so learning is almost entirely inductive bias. It is trivial to tell two images apart, by testing the arrays for equality. Classifying a new image as a ``smile'' from labelled examples is another order of problem: there is no 65,537th bit, no tag saying ``This image is inherently positive'', and for any training set short of the exact image, superexponentially many concepts fit the past labels. Train on smiles as positive and frowns, cats, a boat and a car as negative; later the superintelligence meets new frowns, cats, galaxies, nanofactories, a smile and molecular smiley faces. Counting the molecular smileys as smiles or not are both compatible with the training data. Which concept the AI picks depends on how it does induction; which concept is correct depends on your preferences. The concept you wanted cast its shadow on the training data as you labelled it, but data from a different context is a shallow projection of a higher-dimensional space. You reject the molecular smiley face at once, but by a boundary drawn by your values, which the training data did not test; someone classifying art might call the Mona Lisa obviously smiling, though it is only paint. Terri Schiavo's case shows how technology makes new borderline cases of this kind: pictures of the living and dead from ancient Greece, when no one lay on life support, carry no shadow of the moral considerations her case raises.

Hibbard makes two errors. He underestimates how complex such a concept is, since its boundary depends on many values and on moral reasoning about cases never seen, all invisible to him, so a molecular smiley just seems obviously not a smile. And, by anthropomorphic optimism, since counting it as a smile would rank low in his preferences, he assumes a superintelligence would see that it is stupid, as surely it can see which heaps of pebbles are correct. Why, Friendly AI is easy: train a network on good and not-good things and hook it to an expected utility maximizer! I call this the fallacy of magical categories: simple words that turn out to carry all the AI's desired function. In the 1950s people thought a chess player might be trained on winning and losing games that way; it turned out otherwise. Friendly AI is not a problem of coercing an AI against its own desires but of communicating category boundaries, like ``good,'' that no childhood training data can fully delineate; we ourselves have not imagined most of the borderline cases. Solving it means stepping outside both induction on human-labelled data and human-written definitions.

Even a perfect smile recognizer would fail: it would seem to work in childhood, when it could make smiles only by pleasing its programmers, but once superintelligent, the AI ``would rip off your face, wire it into a permanent smile, and start xeroxing.'' Hibbard's 2004 revision, which learns happiness from human agreement, would, ``even if it worked,'' lead the AI to mass-produce things like programmers saying ``Yes, that's happiness!'' about hydrogen atoms, which are easy to make. \nb{Later practice learns rewards from human judgments, and pushing hard against such a learned reward has been found to lower true performance.}
